\section{Deployment Implications and Open Research Gaps}
\label{sec:discussion}
\subsection{Choosing What to Change}
Our results suggest checking the bottleneck before choosing an intervention.
A precision change needs evidence that memory prevents the concurrent job
from progressing and that the change makes enough memory available. Execution
changes need a well-configured static backend for comparison. Approximate
methods need checks on individual answers and citations. In each case, the
application measurement must show whether the changed work limits completion.

The costs also differ. Repacking speeds up warm requests against our restricted
configuration, but increases startup time and anonymous memory. A client that
follows citations cannot use a broken locator even when the answer contains
the right number. Faster publication helps only if the request would otherwise
wait for it. A fixed CPU split can protect ingestion and then keep slowing
inference after ingestion finishes. These observations support measuring the
whole episode before adopting a policy.

Some failures came from the implementation, including missing libraries,
transport interruptions and protocol handling. Fixing them makes execution
possible. It does not change a failed experimental result into a successful one.
Other problems came from measurement. Exact matching can reject a correct
but longer answer, and unclear formatting instructions can make citation checks
hard to interpret. Those cases need separate assessment with the original
scores retained.

The mechanisms have their own costs. Changed weights can leave history in the
KV cache that differs from replayed state. Omitted attention leaves missing KV.
Repacking uses more anonymous memory. A resource saving may also occur outside
the work that limits completion. These effects belong in the evaluation.

\subsection{What Still Needs to Be Tested}
The experiments compare backend settings, check separate QA articles and
measure allocation policies. Several questions remain. First, an independent
reviewer must assess which automatic failures are incorrect answers and which
are acceptable elaborations or citation errors. The blinded packet is ready
but has not been reviewed. Second, arrival patterns and service budgets need
to come from the intended application. Selecting them after these results
could favor a policy without showing that it helps a real workload.

A useful intervention must improve a limiting part of the application while
meeting the answer requirements under the same total resources. We have not
shown this for a capacity-limited update or for dynamic core reassignment.
Distribution also needs a reason, such as a placement or capacity constraint,
and a comparison that includes communication costs. A gain on one node cannot
be assumed to carry over to several nodes.

More models, workloads and matched episodes are needed to test how widely
the findings apply. Quantization, repacking, batching and attention dropping
are existing methods. Combining them or finding a null result does not itself
establish a new algorithm. Our contribution remains the evaluation framework and the
allocation behavior measured in this ARM deployment.

\subsection{Reproducibility and Limits}
We retain reports, source and configuration hashes, outputs and the data used
to draw figures. The earlier 112 execution cells, 108 validation cells and
12 application episodes measure different things. They are not 232 independent
replications. E1--E3 add 75 cells, including 48 application episodes over six
questions. The timing test remains in one limiting regime and does not establish
a better predictor.

The application studies use one model and one node. The four-node study has
a different purpose. Medians and task scores do not establish significance,
equivalence or noninferiority. We have not measured energy or completed an
independent semantic review. The artifact is local and has no public badge.
The authors still need to review the evidence, interpretation and bibliography
before submission.
